\section{Appchain Cost Model}

\label{sec:costs}
%%==========================================================================

\subsection{Cost Components}

The total cost of operating a Lux appchain has three components:

\begin{enumerate}[leftmargin=1.5em]
  \item \textbf{Infrastructure costs} $C_{\text{infra}}$: Hardware and bandwidth for
    validator nodes.
  \item \textbf{Validation incentives} $C_{\text{val}}$: Rewards paid to validators
    above Primary Network opportunity cost.
  \item \textbf{Protocol fees} $C_{\text{protocol}}$: LUX fees paid to the Primary Network
    (via fee sharing, Section~\ref{sec:feesharing}).
\end{enumerate}

\[
  C_{\text{total}} = C_{\text{infra}} + C_{\text{val}} + C_{\text{protocol}}.
\]

\subsection{Infrastructure Cost Model}

Validator hardware costs are a function of throughput tier. We define five tiers:

\begin{table}[H]
\centering
\caption{Appchain Throughput Tiers and Infrastructure Costs}
\label{tab:tiers}
\begin{tabular}{llrrrrr}
\toprule
\textbf{Tier} & \textbf{Name} & \textbf{TPS} & \textbf{Block Gas} & \textbf{State DB} & \textbf{Node Cost/mo} & \textbf{Min Validators} \\
\midrule
T1 & Micro & $<$100 & 8M & 50 GB & \$120 & 5 \\
T2 & Standard & 100--1K & 30M & 200 GB & \$280 & 7 \\
T3 & High & 1K--10K & 100M & 1 TB & \$650 & 10 \\
T4 & Ultra & 10K--100K & 500M & 5 TB & \$1,800 & 15 \\
T5 & Hyper & $>$100K & 2B & 20 TB & \$4,500 & 20 \\
\bottomrule
\end{tabular}
\end{table}

Annual infrastructure cost for a Tier $k$ appchain with $n_k$ validators:
\[
  C_{\text{infra}}(k) = 12 \cdot n_k \cdot \text{NodeCost}_k.
\]

\subsection{Validation Incentive Cost}

Let $r_{\text{primary}}$ be the annualized Primary Network staking yield and $\delta_k$
be the premium required to attract validators to a Tier $k$ appchain. The annual
validation incentive cost is:
\[
  C_{\text{val}}(k) = n_k \cdot \bar{w} \cdot P_{\text{LUX}} \cdot (r_{\text{primary}} + \delta_k),
\]
where $\bar{w}$ is the mean validator stake weight. We estimate $\delta_k$ empirically
in Section~\ref{sec:validators}.

\subsection{Break-Even Analysis}

For a appchain to be economically sustainable, its fee revenue must exceed total costs:
\[
  R_{\text{fee}} \geq C_{\text{total}}.
\]

Define the \emph{break-even transaction count} as the number of transactions $N^*$
such that average fee per transaction $f$ satisfies $N^* \cdot f = C_{\text{total}}$.
Table~\ref{tab:breakeven} shows break-even estimates by tier:

\begin{table}[H]
\centering
\caption{Appchain Break-Even Transaction Counts (Annual, Assuming \$0.01 Avg Fee)}
\label{tab:breakeven}
\begin{tabular}{lrrrr}
\toprule
\textbf{Tier} & \textbf{$C_\text{infra}$/yr} & \textbf{$C_\text{val}$/yr} & \textbf{$C_\text{total}$/yr} & \textbf{Break-Even Tx} \\
\midrule
T1 & \$7,200 & \$112,000 & \$119,200 & 11.9M \\
T2 & \$23,520 & \$313,600 & \$337,120 & 33.7M \\
T3 & \$78,000 & \$896,000 & \$974,000 & 97.4M \\
T4 & \$324,000 & \$2,688,000 & \$3,012,000 & 301.2M \\
T5 & \$1,080,000 & \$8,960,000 & \$10,040,000 & 1.0B \\
\bottomrule
\end{tabular}
\end{table}

These figures assume LUX at \$8.00, 7\% primary staking yield, and 5\% validator premium
($\delta = 0.05$). They serve as minimum viability benchmarks.

%%==========================================================================
